\section{Introduction}\label{sec:introduction}

Fix $p\in(0,1)$. Let $A_n^{(p)}$ be the real symmetric $n\times n$
matrix whose entries on and above the diagonal are independent and satisfy
\[
 \Prob(A_{ij}^{(p)}=1)=p,\qquad
 \Prob(A_{ij}^{(p)}=-1)=1-p \qquad (i\le j).
\]
The entries below the diagonal are determined by symmetry. We study the
exponential rate at which this matrix is singular.

\begin{theorem}\label{thm:main}
For every fixed $p\in(0,1)\setminus\{1/2\}$,
\[
 \lim_{n\to\infty}\Prob\bigl(\det A_n^{(p)}=0\bigr)^{1/n}
   =p^2+(1-p)^2.
\]
\end{theorem}

The quantity on the right is the probability that two independent signs
with this bias agree. Requiring two rows of the matrix to agree gives the
same exponential rate: symmetry changes only the probability of the fixed
corner where those rows meet. The upper bound shows that all linear
dependencies together have this rate. The assertion concerns each fixed
$p$; its error term may depend on $p$.

\subsection{Historical context}

The singularity problem for discrete random matrices goes back to
Koml\'os, who proved asymptotic nonsingularity for matrices with
independent uniform \(0/1\) entries~\cite{Komlos}.
For independent uniform signs, Kahn, Koml\'os, and Szemer\'edi obtained
an exponential upper bound~\cite{KahnKomlosSzemeredi}, and Tikhomirov
identified the sharp exponential rate \((1/2+o(1))^n\)~\cite{Tikhomirov}.
Jain, Sah, and Sawhney determined the singularity asymptotics for every
fixed nonuniform finitely supported entry law~\cite{JainSahSawhneyDiscrete}.
In particular, their result gives the exponential rate
\(p^2+(1-p)^2\) for independent biased signs.
A recent preprint of Hu, Song, and Wu establishes the full asymptotic
\((2+o(1))n^2 2^{-n}\) for independent uniform signs~\cite{HuSongWu}.
All these models have independent rows; symmetry removes that independence.

Costello, Tao, and Vu proved asymptotic nonsingularity for symmetric
Bernoulli matrices by developing quadratic Littlewood--Offord
estimates~\cite{CTV}. Their general maximum-atom theorem also covers
each fixed biased sign law. In the unbiased symmetric sign model,
Nguyen obtained arbitrary polynomial decay~\cite{Nguyen}, Vershynin
obtained a stretched exponential bound~\cite{Vershynin}, and further
combinatorial progress included the work of Ferber and
Jain~\cite{FerberJain}. Campos, Jenssen, Michelen, and Sahasrabudhe
proved the exponential bound \(\exp(-cn)\)~\cite{SymmetricExponential}.
The companion manuscript~\cite{SymmetricRate} proves the sharp unbiased
rate \((1/2+o(1))^n\). The present paper establishes its fixed-bias
analogue, at the level of the exponential rate.

\subsection{The proof and its main ingredients}

The main difficulty is to combine symmetry with the unequal probabilities
of individual columns. Multiplying the whole matrix by $-1$ allows us to
assume $p<1/2$. A prescribed column can then have probability as large as
$(1-p)^{n-1}$. That bound alone is too large for the rate in
Theorem~\ref{thm:main}. We obtain two different gains: a structural
exception requires two rare column extensions, and a nonexceptional
singular matrix admits a deletion whose column has nearly its full
Shannon entropy.

Section~\ref{sec:reductions} first reduces to corank one. For a nonzero
kernel vector, choose a coordinate of largest absolute value as its anchor and normalize
the remaining coordinates, called its tail, to have Euclidean norm one.
Tails close to sparse vectors are controlled directly at the claimed
rate. For the remaining vectors, we multiply by a global sign and change
coordinate signs to make the anchor row all ones. Its Schur complement
is $-2$ times a symmetric matrix with entries in $\bits$.
Conditional on the anchor row, these entries are
independent on and above the diagonal, their probabilities belong to
$\{p,1-p\}$, and their mean matrix has rank at most two. Deleting a
suitable coordinate gives a nonsingular principal minor $B$. The omitted
column $z$ must satisfy
\[
 z^{\mathsf T}B^{-1}z\in\bits,
 \qquad \lVert B^{-1}z\rVert_2\le L_0\sqrt n
\]
for a fixed constant $L_0$. The coordinate is chosen so that the original
sign column also has its typical number of positive entries. This use
of the gap between entropy and collision probability follows the
independent-entry strategy of Jain, Sah, and Sawhney~\cite[Section~6]{JainSahSawhneyDiscrete}. Thus each
compatible column has probability governed by the binary entropy of $p$.
Proposition~\ref{prop:column-bound} is the central remaining assertion:
there are few such columns except on an event with a probability close
to $(1-p)^{2n}$. Both bounds leave a strict margin beyond the required
rate when $p\ne1/2$.

The arithmetic part follows the principal-minor and discriminant-group
method of the companion paper~\cite{SymmetricRate}. We use its
low-height subspace cover, robust lattice bounds, and deterministic
pairing and minor identities; the precise statements needed are given
where they enter the proof. We strengthen the companion's cube-intersection estimate to accommodate the biased law and to obtain
nearly complete loss of cube entropy, simultaneously over intermediate
lattices. For a nonsingular bit matrix
$B$ of size $m$ close to a reference dimension $N$, consider
\[
 B\Z^m\subseteq K\subseteq\Z^m.
\]
Theorem~\ref{thm:small-intersection} shows that, outside an event of
probability smaller than every fixed exponential in $N$, an index as
large as $N^{\tau N}$ forces every coset of $K$ to meet the cube in
fewer than $2^{\alpha N}$ points. Here $\alpha>0$ can be any fixed
constant, and $\tau$ may decrease as a fixed inverse power of
$\log\log N$. The simultaneous conclusion permits the lattice to be
chosen after the matrix has been observed.

The proof of this estimate, in Section~\ref{sec:lattice}, converts a
large cube intersection into many effective linear tests from a dual
lattice. A low-height cover, combined with a change of scale, confines
those tests near a subspace of
dimension strictly less than $N$. Their approximate locations can
therefore be counted more cheaply than independent groups of entries can satisfy
their integrality conditions. Gaussian sampling on the dual lattice
makes this comparison effective for the biased distribution.

Sections~\ref{sec:pairings} and~\ref{sec:paths} use the estimate in
two ways. First, the finite pairing on $\Z^m/B\Z^m$ bounds the
number of admissible columns when all but the two largest invariant
factors of the group have small total logarithmic size. Second, paths
of nonsingular principal minors control the
remaining matrices. In a short terminal portion of a path, a matrix
with many admissible columns forces either two rare extensions or one
extension by two rare columns. Both alternatives have the required
probability cost. This proves Proposition~\ref{prop:column-bound}
and completes the reduction to Theorem~\ref{thm:main}.
